\subsection{The original local chart and integral Selmer comparison}
\label{sec:support-local-chart}

We now construct the local chart at the diagonal representation.  The
chart is over the fraction field.  The local and global equation ideals
used in the integral length comparison remain the original ideals.

Let $O$ be the integers of a finite extension $E/\mathbf Q_\ell$,
where $\ell\ne2$.  On $G=G_{\mathbf Q_\ell}$ retain the original
unramified quadratic character $\eta$, finite character $\tau$, and
continuous character $\Psi$ satisfying
\[
 \Psi|_{I_\ell}=\varepsilon^e\tau|_{I_\ell},\qquad e\ge1.
                                                               \tag{LC1}
\]
Set $G_L=\ker\eta\cap\ker\tau$.  This fixes the finite local
extension $L$; its degree is not bounded in this argument.  Let $R$ be
the universal framed local ring for the residual representation of
\[
 r_0(g)=\bigl(\operatorname{diag}(\Psi(g),1,\Psi(g)^{-1}),
                    \eta(g)\bigr)
 \quad\hbox{in}\quad
 \mathrm{GL}_3\rtimes\{\pm1\},
                                                               \tag{LC2}
\]
where the nonidentity component acts by $A\mapsto JA^{-\mathsf t}J$
and $J$ reverses the basis.  Write $x:R\to O$ for its base point.
The multiplier is fixed to be $\eta$.

Use exactly the ordinary ideal $J_{\rm ord}$ and element $\delta$
of Proposition~\ref{prop:selmer-crystalline-components}.  Thus
$J_{\rm ord}$ imposes the characteristic polynomials and the ordered
three-factor products on $G_L\cap I_\ell$, and
$\delta=\ell d_0d_1$ uses the original inertia element and Frobenius
in $G_L$.  Put
\[
 J_Y=J_{\rm ord}:\delta^\infty,\qquad
 J_{\rm red}=\sqrt{J_Y},\qquad
 J_t=\bigcap_{\substack{P\text{ minimal over }J_{\rm red}\\
                               P\subset\ker x}}P.
                                                               \tag{LC3}
\]
The universal property makes the base point $x$ unique, so the last
condition is the source's condition at every base point.  The selected
local ring is $R_\ell=R/J_t$.

\begin{theorem}[The actual local component and chart]
\label{thm:support-local-chart}
There is exactly one prime in the intersection in (LC3), and it is
$J_t$.  The ring $R_\ell$ is stable under conjugation by every matrix
congruent to one over every complete local coefficient algebra.
There are twelve actual matrix-entry differences $u_i\in\ker x$
and a homomorphism
\[
 \phi:R\longrightarrow E[[z_1,\ldots,z_{12}]],\qquad
 \phi(u_i)=z_i,\qquad \ker\phi=J_t,
                                                               \tag{LC4}
\]
whose constant term is $x$.  The entries are homogeneous for
conjugation by $\operatorname{diag}(t,1,t^{-1})$, with the original
matrix weights.  The coordinate multiplicities are
\[
\begin{array}{c|rrrrr}
 j&-2&-1&0&1&2\\ \hline
 \#\{i:\deg z_i=j\}&1&2&3&4&2.
\end{array}                                                    \tag{LC5}
\]
If $x(a)=0$ and the linear part of $\phi(a)$ vanishes, some
$0\ne c\in O$ satisfies
$ca\in (\ker x)^2+J_t$.  Every linear coordinate direction is
represented by an integral element of $R$, up to a nonzero scalar
in $O$.  These are the full original local-chart requirements
\cite{support-local-chart-source}.
\end{theorem}

\begin{proof}
We first identify the completed characteristic-zero local ring, then
return to the integral selected component.

\emph{The crystalline deformation ring.}
Fix the Hodge type and inertial type of $r_0$ in (LC2).  They are
defined over $E$; on $G_L$ the type is crystalline with the three
distinct Hodge--Tate weights $-e,0,e$, using the convention that
$\varepsilon$ has weight $-1$.  The characteristic-zero framed
deformation ring of this type is regular.  We use
Bellovin--Gee, Theorem~3.3.7, which allows disconnected reductive
groups \cite{support-local-chart-BG}.  Here the group is
$\mathrm{GL}_3\rtimes\{\pm1\}$; its identity component has
dimension nine and its regular Hodge flag has dimension three.
The component character is discrete and fixed by reduction.  Its
centre is smooth because $\ell$ is odd.  Thus the completion at
$r_0$, with its specified residue-field map to $E$, is
\[
                         C\simeq E[[z_1,\ldots,z_{12}]].       \tag{LC6}
\]
Only this complete local regularity is used, not a claim of integral
smoothness.  The cited theorem follows by imposing zero monodromy
in the crystalline deformation problem: the obstruction complex has
two terms in degrees zero and one.  This also explains why a possible
equality of adjacent Frobenius ratios with the cyclotomic character
does not invalidate (LC6).

\emph{Crystalline deformations near the base are ordinary.}
We include this step because smoothness of an unrelated local ring
would not prove the theorem.  Let $B$ be a finite local Artinian
$E$-algebra, and let a deformation of $r_0$ to $B$ have the fixed
potentially crystalline type.  Restrict to $G_L$ and take its filtered
crystalline module.  After a finite scalar extension, a positive
iterate of Frobenius acts linearly on each coefficient factor.  Its
three residual eigenvalues have distinct valuations, proportional to
$-e,0,e$.  Hensel factorization of its characteristic polynomial
therefore gives three rank-one generalized eigenspaces over $B$.
The sums of the first one and first two, ordered by increasing
valuation, are independent of the iterate and the splitting extension.
They are stable under Frobenius and the finite descent action.

Write $D_{\le k}$ for the first $k$ slope lines.  At every embedding
of $L$, the Hodge filtration of the split base is opposite to this
slope flag.  In particular $D_{\le1}$ misses the two-dimensional
filtration step, $D_{\le2}$ misses the one-dimensional step, and
$D_{\le2}$ meets the two-dimensional step in dimension one.
These are the minimal possible intersection dimensions.  The minors
expressing these transversality conditions are nonzero at the base,
hence units over $B$.  Consequently all intersections with the Hodge
filtration, and all quotient filtrations, are direct summands of the
same ranks over $B$.

This also verifies strictness for the nilpotent coefficient
filtration, a point needed before applying field-valued comparison
theorems.  Choose a $B$-basis of each $D_{\le k}$ adapted to its
induced Hodge filtration.  For every $n$ and every filtration step
$F$, one then has
\[
 (\mathfrak m_B^nD_{\le k})\cap F
                  =\mathfrak m_B^n(D_{\le k}\cap F).
\]
The associated filtered Frobenius modules for the
$\mathfrak m_B$-adic filtration are therefore
$(\mathfrak m_B^n/\mathfrak m_B^{n+1})\otimes_E D_{\le k,0}$;
the same assertion holds for successive slope quotients.  They are
sums of the corresponding admissible modules of the split base.
Weak admissibility is preserved by strict extensions: Newton and
Hodge numbers are additive, and the required inequalities for a
subobject follow by intersecting it with the subobject and projecting
it to the quotient.  Induction through the finite coefficient
filtration proves that every $D_{\le k}$ and successive quotient
is weakly admissible.

Apply the Colmez--Fontaine equivalence to these finite-dimensional
$\mathbf Q_\ell$-objects \cite{support-local-chart-CF}.  Full
faithfulness retains the $B$-action, and exactness carries the
strict coefficient filtrations to the corresponding representation
filtrations.  The representation associated with $D_{\le k}$ has
reduction of $E$-dimension $k$ and total $E$-dimension
$k\dim_E B$.  Nakayama gives a surjection from $B^k$, which is an
isomorphism by those dimensions.  Thus we obtain a free
rank-one/rank-two Galois flag with free successive quotients.
Uniqueness of the Frobenius flag makes the construction descend
through the scalar extension and respect the polarization.

On each rank-one quotient, fixed inertial type and fixed Hodge weight
make the inertia character exactly $\Psi,1,\Psi^{-1}$, in that
order.  There is no infinitesimal variation of its finite inertia
character: averaging over the finite inertia quotient is valid in
characteristic zero.  Nor can its Hodge weight vary in a crystalline
Artinian deformation.  Consequently every such deformation satisfies
the literal equations of $J_{\rm ord}$.  This proves that their
images in the reduced ring $C$ vanish.

\emph{The saturation selects this same component.}
Proposition~\ref{prop:selmer-crystalline-components} proves the
opposite inclusion: every finite-algebra point of $R/J_{\rm red}$
is crystalline on the same $G_L$.  Near $x$ its finite inertial
type and Hodge type are those of $r_0$.  Indeed the descent inertia
factors through a fixed finite group.  Its representations over
Artinian $E$-algebras are rigid up to conjugacy: averaging kills
the first cohomology at each square-zero coefficient step.  The
Hodge ranks are locally constant ranks of the filtered direct
summands and reduce to the ranks at $r_0$.  Thus it factors through the
fixed-type crystalline deformation ring, including at nilpotent
Artinian points.  In the other direction, $\delta$ is a nonzero
element of $C$.  To check this even when $\delta(x)=0$, vary
$\Psi$ by an unramified character $\nu$ with Frobenius value
$1+T$.  The diagonal representation (LC2) with $\Psi\nu$ has the
same crystalline type.  If the original Frobenius is of degree
$m\ge1$ in the unramified quotient, both adjacent factors become,
up to nonzero constant factors,
\[
                  \Psi(\mathrm{fr})(1+T)^m
                           -\varepsilon(\mathrm{fr}).         \tag{LC7}
\]
This power series is nonzero in characteristic zero.  The original
inertia eigenvalue differences are also nonzero, since the selected
cyclotomic inertia value has infinite order.  Hence $\delta$ is
nonzero in the domain $C$.

Since $J_{\rm ord}$ vanishes in $C$ and $\delta$ is not a zero
divisor there, its saturation and radical vanish there as well.
Conversely, the preceding finite-algebra comparison implies that the
completed local quotient by $J_{\rm red}$ factors through $C$.
These two maps are inverse: both are quotients of the completed
universal local ring and agree on its universal representation.
The completion of $(R/J_{\rm red})[1/\ell]$ at $x$ is therefore
exactly $C$.

Here the local passage is legitimate for the full ideals.  The
universal ring is a quotient of a finite power-series ring over $O$
and is excellent.  Localization and completion are flat; completion
of its reduced local quotient is reduced.  The closed-locus comparison
uses Kisin's characterization on every finite coefficient algebra,
as proved in Proposition~\ref{prop:selmer-crystalline-components}.
It is not an inference that field points detect nilpotents.  The
domain $C$ and faithful flatness of local completion show that
exactly one minimal prime of $J_{\rm red}$ passes through $x$.
Call it $P$.  Its contraction is the original $J_t=P$; thus
$R_\ell$ is a domain, and its map to $C$ is injective.

\emph{The weights and integral entry coordinates.}
The tangent space of the framed crystalline functor is the space
$Z_f^1(G,\operatorname{ad}r_0)$ of cocycles whose classes satisfy the
finite local condition.  The weight-$j$ adjoint subspace has matrix
dimension $d_j$, where
\[
                    (d_{-2},d_{-1},d_0,d_1,d_2)=(1,2,3,2,1).
\]
The involution on matrices is $X\mapsto-JX^{\mathsf t}J$.
It preserves each weight.  Since $2$ is invertible, its reflected
pairs split into their two eigenlines.  Each resulting character
has Tate exponent $ej$, the finite twist $\tau^j$, an unramified
factor, and possibly the unramified quadratic factor $\eta$.
Bloch--Kato gives zero finite first cohomology for negative exponent
and one dimension for positive exponent.  Adding coboundaries gives
one cocycle dimension per negative matrix direction and two per
positive direction.  At weight zero the characters are unramified;
the finite classes are unramified classes, whose dimension plus the
coboundary dimension is one per direction.  This gives
\[
 \dim_E Z_f^1(G,\operatorname{ad}r_0)_j
       =\begin{cases}d_j,&j\le0,\\2d_j,&j>0,\end{cases}
                                                               \tag{LC8}
\]
and hence (LC5).  The same calculation applies with the finite
character in (LC1): its inertia twist does not change the nonzero
Tate-exponent formula, and the weight-zero factors remain unramified.

Actual matrix-entry differences separate framed tangent vectors.
Choose twelve of their linear parts forming a basis, weight by
weight.  The formal inverse-function theorem makes these entries
parameters in $C$.  Their transformation under the diagonal torus
is exactly by their matrix weights, so the parameter isomorphism and
all matrix entries are homogeneous.  The composite $R\to C$ is
the map (LC4), with the asserted kernel.

To check the integral tangent assertion, put $\mathfrak p=\ker x$.
If $x(a)=0$ and the linear part of $\phi(a)$ vanishes, faithful
flatness gives $a\in\mathfrak p^2+J_t$ after localization at the
characteristic-zero point.  Clear denominators to obtain
$ba\in\mathfrak p^2+J_t$ for some $b\in R$ with $x(b)\ne0$.
Since $(b-x(b))a\in\mathfrak p^2$, the required scalar is
$c=x(b)$.  Clearing an additional power of $\ell$, if needed,
keeps $c\ne0$.  The same localization argument supplies finite
jets up to a nonzero scalar in $O$; in degree one the selected
integral entries already supply every coordinate.

\emph{Conjugation on the integral component.}
In $R[[Z_{ab}]]$, conjugation by the universal matrix $1+Z$ is an
automorphism with inverse given by $(1+Z)^{-1}$.  Characteristic
polynomials, the ordered product equations, and $\delta$ are
conjugation invariant, so this automorphism preserves $J_{\rm red}$
extended to the power-series ring.  It fixes the base point at $Z=0$.
The unique minimal prime through that point is $P[[Z]]$, since its
completed local quotient is $C[[Z]]$, a domain.  Thus it preserves
$P[[Z]]$.  Specializing $Z$ to the entries of any matrix in
$1+\mathfrak m_A M_3(A)$ proves the stated conjugation stability
over every complete local algebra $A$, including algebras with
nilpotents and $\ell$-torsion.
\end{proof}

\begin{theorem}[The actual integral local lifting]
\label{thm:support-actual-integral-lifting}
For each original good global datum, the actual marked global model
exists with the original local ring $R_\ell$ and full global
equation ideal.  Every original normalized strict cocycle modulo a
nonzero $a\in O$ gives a weight-$-1$ derivative on that model which
lifts to an integral derivative on its ambient algebra.  The modulus
remains $a$.  In particular the strict Selmer module embeds into
the dual of the original negative conormal quotient.
\end{theorem}

\begin{proof}
Choose a finite weighted presentation of $R_\ell$ by actual
matrix-entry differences.  Such differences topologically generate
the framed ring, and the preceding conjugation action gives their
matrix weights and a homogeneous kernel.  Theorem
\ref{thm:support-relative-presentation} supplies the relative global
presentation over this very local presentation.  Before imposing the
global and marking equations, its ambient complete algebra is
\[
                       \widehat A=R_\ell[[X_{j,s}]].           \tag{LC9}
\]
Let $A$ be its algebra of finite sums of homogeneous elements, and
let $Q$ be the full homogeneous ideal of global and marking equations
in $A$.  With $a_j$ as in that theorem, put
\[
 r=a_{-1}+2,\quad s=a_{-2}+1,\quad
 r'=a_1+3,\quad s'=a_2+1,\quad m=a_0+3.
\]
The numbers of equations of weights $-1,-2,1,2,0$ are
$r+1,s,r',s',m$.  The local chart and the free variables give
chart dimensions $r,s,r'+1,s'+1,m$ in those same weights.
All entries have homogeneous preimages in $A$.  Its coordinate
derivatives are integral multiples of derivatives of entries of the
same weight; inverse-matrix entries have this property by
$d(B^{-1})=-B^{-1}(dB)B^{-1}$ at the diagonal base.

For completeness, the passage from algebraic points of $A/Q$ to
local points of the completed marked ring adds no arithmetic input.
Its weight-zero quotient is finite and local over $O$ by
Theorem~\ref{thm:support-actual-integral-finiteness}, with the
original augmentation as its residual point.  For its maximal ideal
$\mathfrak n$, some power of $\mathfrak n$ is contained in
$\ell$ times that quotient.  Thus every characteristic-zero point
sends $\mathfrak n$ to elements of strictly positive valuation.
Opposite-weight generator products of weight zero belong to
$\mathfrak n$, as do the weight-zero augmentation generators.
Consequently the strict
inequalities making all nonzero-weight generators small after
scaling by $t^j$ have a common solution for $v(t)$: each lower
bound is smaller than each upper bound precisely by one of those
product inequalities.  Weight-zero generators are already small.
Choose such $t$ in a finite coefficient extension.  Power-series
evaluation then gives the local point.  The same convergent
first-order substitution extends any derivative annihilating $Q$;
its finitely many generator values lie in one finite coefficient
field.  This proves the point and derivative parts of the original
model, as well as its chart and equation counts.

Now take a normalized strict cocycle $\gamma$ with coefficients
$O/a$ in weight $-1$.  On the local group it is
$X-\operatorname{Ad}(r_0)X$ for a weight-$-1$ matrix $X$ over
$O/a$.  Lift its two possibly nonzero entries to a matrix $Y$ over
$O$.  Both $\eta$ and $\Psi$ are continuous, so conjugating the
integral base representation by $1+\epsilon Y$ over
$O[\epsilon]/(\epsilon^2)$ gives a continuous integral local
deformation.  Theorem~\ref{thm:support-local-chart} puts it on
the original $R_\ell$.  Its derivative is integral and reduces
to the required local coboundary.

The square-zero deformation defined by $\gamma$ satisfies the
actual global equations: its normalization gives the two markings,
it vanishes on the prescribed killed inertia subgroups, and its
local restriction is the just constructed point of $R_\ell$.
Universality gives a derivative $D$ on the original global ring.
Extend its integral local derivative to (LC9) by lifting the values
of $D$ on the new variables of weight $-1$ and using zero at every
other weight.  On a power series this derivative reads only its
linear coefficients.  It kills the extended local ideal and
reduces to $D$, because both agree on the entry coordinates.
This proves the original integral lifting assertion.

Write $I=\ker(A\to O)$ and
$I^{(2)}=\{f:\exists,0\ne c\in O, cf\in I^2\}$.
The integral derivative kills $I^{(2)}$ because its target $O$ is
torsion-free, and its reduction kills $Q$.  The established normalized
cocycle construction therefore gives the compatible injections
\[
 \mathcal S[a]\hookrightarrow
 \operatorname{Hom}_O\!\left(
    (I/(I^{(2)}+Q))_{-1},\,O/a\right),\qquad
 \mathcal S\hookrightarrow
 \operatorname{Hom}_O\!\left(
    (I/(I^{(2)}+Q))_{-1},\,E/O\right).                         \tag{LC10}
\]
Here $\mathcal S$ is the original strict module with character
$\Psi^{-1}$.  The original local unit condition makes a local
primitive of an $a$-torsion class itself $a$-torsion, so no modulus
is enlarged.  Compatibility follows from the coefficient map
$O/a\to O/(ab)$, $u\mapsto bu$.  Matrix entries recover
$\gamma$, proving injectivity.  No torsion is removed after $Q$
is imposed; that torsion is the quantity measured in (LC10).
\end{proof}

\begin{corollary}[The original one-character trace inequality]
\label{cor:support-one-character-trace}
Retain every original good-data hypothesis for the character $\chi$,
twist $e$, imaginary quadratic field $K$, and coefficient prime
$\ell$.  For $\sigma\in G_K$ put
$t_\sigma=\Psi(\sigma)+1+\Psi(\sigma)^{-1}$.
Suppose every original admissible rank-three representation, including
the original rank-two-plus-line alternative, has trace at $\sigma$
different from $t_\sigma$.  There is a finite set $\mathcal X$ of
pairwise nonisomorphic such representations for which
\[
 \frac{\operatorname{length}_O\mathcal S}
      {e(O/\mathbf Z_\ell)}
 \le 2\sum_{\rho\in\mathcal X}
       v_\ell\bigl(\operatorname{tr}\rho(\sigma)-t_\sigma\bigr),
 \qquad v_\ell(\ell)=1.                                     \tag{LC11}
\]
The admissibility bounds are the original $e,N$ bounds in rank
three and $e,N|\operatorname{disc}K|$ bounds in the rank-two
alternative.  In particular $\mathcal S$ has finite length.
\end{corollary}

\begin{proof}
Apply Theorem~\ref{thm:selmer-abstract-length} to the actual model
just constructed.  The negative and zero cotangent spanning,
fixed-point separation, integral weight-zero finiteness, and
two-sided tangent bounds are the previously proved arithmetic
statements, including Theorems~\ref{thm:support-type-one-rigidity}
and~\ref{thm:support-type-two-rigidity}.  Use as trace element the
actual difference at $\sigma$.  The nonvanishing assumption is
exactly the hypothesis required by that theorem.  Combining its
upper length bound with (LC10) gives (LC11), before division by
the coefficient ramification index.

The actual type-I and type-II identifications put these very
representations in the stated admissible families.  We verify the
isomorphism issue, since otherwise the trace sum could acquire an
unjustified multiplicity.  Write $T=\operatorname{diag}(1,0,-1)$
and $g^*=cg^{-1}c$.  The original marking gives
$r(g_1)-r(g_1^*)=\zeta T$, with $\zeta$ congruent to a fixed
unit at every local point.  For two such points, $\zeta$ and
$\zeta'$ are nonzero and $\zeta'\ne-\zeta$: their common
residue is a unit and $2$ is a unit.  An intertwiner $M$ on $G_K$
satisfies $M\zeta T=\zeta'TM$.  The three diagonal eigenvalues
and the preceding exclusion force every off-diagonal entry of $M$
to be zero.  Write $M=\operatorname{diag}(m_0,m_1,m_2)$, with
all $m_i\ne0$.

For a matrix $A=r(g)$ and $A'=r'(g)$, $g\in G_K$, the
intertwining identities at $g$ and $g^*$ give
\[
 m_i A_{ij}=A'_{ij}m_j,\qquad
 m_{2-j}A_{ij}=A'_{ij}m_{2-i}.
\]
If $m_0m_2=m_1^2$, put $t=m_0/m_1$; these identities give
$A'_{ij}=t^{j-i}A_{ij}$.  Otherwise they force all entries
joining the middle coordinate to either outside coordinate to
vanish.  Choose $t^2=m_0/m_2$ in the algebraically closed
coefficient field; the same equality then holds for the remaining
entries.  Hence $r'$ is conjugate to $r$ by
$\operatorname{diag}(t,1,t^{-1})$ on $G_K$.  Both markings at
$c$ equal the identity matrix and are preserved by this torus, so
the equality holds on the full group.  Thus the original marking
identifies isomorphic two-sided representations precisely up to
the allowed torus action.
Take one representative from each orbit in the finite cover of the
abstract length theorem.  This gives pairwise nonisomorphic
representations and preserves their traces; no extra multiplicity
factor is introduced \cite{support-local-chart-source}.

For a general original coefficient ring, the prime-to-$\ell$
character values lie in a finite unramified coefficient subring
$O_0$.  Flat extension of the same local/global cochain kernels
gives $\mathcal S_O=\mathcal S_{O_0}\otimes_{O_0}O$.
Length is multiplied by $e(O/O_0)$, and dividing by
$e(O/\mathbf Z_\ell)$ cancels this factor.  The trace values in
$\overline{\mathbf Q}_\ell$ are unchanged.  This supplies the
full original coefficient quantifier.
\end{proof}

This closes the local-chart and integral-comparison requirements for
one good character.  It does not bound the sum in (LC11) throughout
the growing family.  The next arithmetic task is to realize these
same admissible trace families at the auxiliary Frobenius elements,
with their conductor, cardinality and integral polynomial bounds,
and to sum the resulting valuations with the original weights and
exceptional primes.  No uniform Selmer product estimate follows
from the fixed-data inequality alone.
